\section{Rational reachability and the signed core}\label{sec:spectral}

Let \(J\) be a finite set of response coordinates. A held test has the form
\begin{equation}\label{eq:binomial}
 x^{u}=c x^{v},\qquad u,v\in\NN^J,\quad c\in\QQ,
 \qquad x^u=\prod_{j\in J}x_j^{u_j}.
\end{equation}
We use \(x_j^0=1\), including at \(x_j=0\). The guards are finitely many
absolute-sum budgets \(\sum_{j\in J_b}|x_j|\le b\), with rational
bounds and coordinate sets \(J_b\subseteq J\). Write
\(\mathcal F_{\QQ}\) for the lawful rational points. Assume that
\(\mathcal F_{\QQ}\) is bounded in \(\RR^J\). Covering \(J\) by
the budget sets suffices, as for the mark budgets in \eqref{eq:guards}.
The real lawful set is closed since the tests are polynomial equalities
and the guards are closed inequalities.
Two lawful rational points are \emph{finely connected} when for every
positive rational \(\varepsilon\) a finite chain in \(\mathcal F_{\QQ}\)
joins them with all successive sup-norm differences less than
\(\varepsilon\).

\subsection{Support reduction}

For a point \(x\), its \emph{signed chart} specifies \(S=\supp x\) and the
sign of \(x_j\) for every \(j\in S\). A row \((u,v)\) is active on \(S\)
when \(\supp u\cup\supp v\subseteq S\).
Let \(A_S\) be the matrix of row differences \(u-v\) for these active tests,
restricted to columns in \(S\).
On an inhabited chart, an active test has nonzero \(c\), compatible signs
and two nonzero monomials. An inactive test with nonzero \(c\) has
both monomials zero. A test with \(c=0\) requires its left monomial to vanish,
and this remains true on every smaller support. In the native class the
zero-coefficient case occurs only in a zero calibration.

\begin{definition}[Contraction and signed core]\label{def:signed-core}
Write \(S\to T\) if there are integer weights \(w_j\geq0\) on \(S\) such that
\[
  A_Sw=0,\qquad T=\{j\in S:w_j=0\}.
\]
A contraction keeps the signs on \(T\) and assigns zero off \(T\).
In a sequence of contractions use the matrix of the current support at
each step. Denote the least support reached from \(S\) by \(\Core(S)\)
and define
\[
 \SCore(x)_j=
 \begin{cases}
   \sgn(x_j), & j\in\Core(\supp x),\\
   0, & j\notin\Core(\supp x).
 \end{cases}
\]
The test list is fixed when constructing this map.
\end{definition}

\begin{lemma}[Realization and restriction]\label{lem:contraction}
A contraction on an inhabited rational signed chart can be realized
from every rational point of that chart by arbitrarily fine rational
chains with the indicated endpoint signs. Moreover, \(S\to T\)
restricts to \(U\to U\cap T\) for every \(U\subseteq S\).
The least reachable support in \Cref{def:signed-core} exists, is unique,
and is unchanged by further reductions.
\end{lemma}
\begin{proof}
For realization, let \(x\) be the initial rational point and set
\(x_j(t)=x_jt^{w_j}\) for \(j\in S\) and \(x_j(t)=0\) otherwise, on
\(0\le t\le1\), again with \(t^0=1\). For \(t>0\), equal total weights
preserve every active test and the inactive tests stay true. At \(t=0\)
the polynomial equalities hold by continuity. Absolute-sum budgets
can only decrease. The resulting curve is continuous, and its samples at
rational parameters are rational, including its limiting endpoint.
A sufficiently fine rational mesh gives the claimed chains.

For restriction, a row active on \(U\) was active on \(S\) and has no
factors outside \(U\). Cutting \(w\) down to \(U\) therefore preserves its
weight balance. The zero-weight coordinates in \(U\) are \(U\cap T\).
Repeated restriction gives the same intersection rule for an entire
reduction path.

Choose a reachable support \(C\) of minimum cardinality. For every reachable
\(U\), restrict the path from \(S\) to \(C\) to \(U\); this reaches
\(U\cap C\) from \(U\). Since \(U\cap C\) is also reachable from \(S\),
minimality forces \(C\subseteq U\). Thus \(C\) is least, is reachable from
every reachable support and is still least there.
\end{proof}

\paragraph{Finite calculation.}
For a prescribed survivor set \(T\), impose \(w_j=0\) on \(T\) and
\(w_j\ge1\) elsewhere on \(S\) in the rational linear system \(A_Sw=0\).
Clearing denominators produces an integer witness from any feasible
rational one. To compute the core, at each support test, for each \(j\),
whether the nonnegative balance system admits \(w_j\ge1\).
The sum of one witness for each feasible \(j\) removes the entire set of
currently removable coordinates in one step. Recompute the matrix and
repeat. Every nonempty round removes a coordinate, and a stable support
has no proper reduct. By \Cref{lem:contraction} it is the core.

\subsection{Limits of charts}\label{sec:chart-limits}

\begin{lemma}[A chart limit gives a contraction]\label{lem:limit}
If \(b\) lies in the real closure of a lawful signed chart on support
\(S\), then the chart contracts to \(\supp b\) with the signs of \(b\).
This holds even when the particular meeting point \(b\) is irrational.
\end{lemma}
\begin{proof}
Take a sequence \(x^{(k)}\) from the chart converging to \(b\).
Put \(T=\supp b\) and \(Z=S\setminus T\).
Surviving coordinates retain their signs by convergence.
If \(Z\) is empty, use zero weights. Otherwise, with
\(\ell^{(k)}=\log|x^{(k)}|\) on \(S\), the active equations are
\[
 A_{S,T}\ell^{(k)}_T+A_{S,Z}\ell^{(k)}_Z=\log|c|,
\]
where the right side lists the log-magnitudes of the active coefficients.
The vector \(\ell^{(k)}_T\) is bounded, while every entry of
\(\ell^{(k)}_Z\) tends to \(-\infty\). Hence
\(A_{S,Z}\ell^{(k)}_Z\) is bounded. Choose a linear right inverse
\(B\) on the image of \(A_{S,Z}\). The vectors
\[
 \ell^{(k)}_Z-B A_{S,Z}\ell^{(k)}_Z\in\ker A_{S,Z}
\]
are eventually strictly negative in every coordinate, because only a
bounded correction was subtracted. Negation gives a strictly positive
vector in the real kernel. The matrix has integer entries, so a rational
basis of its kernel and rational approximation give a positive rational
vector as well. Clear denominators and extend by zero on \(T\).
The resulting integer weights give precisely the required contraction.
\end{proof}

\subsection{Exact rational motion within a chart}\label{sec:rational-chart}

\begin{lemma}[Rational chart interpolation]\label{lem:rational-chart}
Any two lawful rational points of the same signed chart are finely connected
through rational points of that chart.
\end{lemma}
\begin{proof}
Fix endpoints \(x,y\) on support \(S\). Their real log segment has coordinates
\[
 \xi_j(t)=\sgn(x_j)|x_j|^{1-t}|y_j|^t \quad(j\in S)
\]
and is zero outside \(S\). It preserves the active equations, since their
log form is affine. Each absolute-sum budget is a convex function of the
log vector. At \(0<t<1\) a budget containing a changed coordinate has
strict slack, by strict convexity of the exponential in that coordinate.
Every unchanged block has exactly its endpoint values throughout.

To obtain rational samples that satisfy the equalities exactly, factor
each positive rational ratio
\(\rho_j=|y_j|/|x_j|\) as \(\prod_p p^{z_{jp}}\), with integer exponents
and a finite set of primes shared across the coordinates. The endpoints
satisfy the same active equations, so \(\rho^{u-v}=1\) for each active
row. Unique factorization gives \(A_S z_p=0\) for every vector \(z_p\)
of prime exponents. All \(z_{jp}\) vanish where \(\rho_j=1\).
At an interior time \(t\) choose positive rationals \(r_p\) close to \(p^t\).
The coordinates \(x_j\prod_p r_p^{z_{jp}}\) are rational, have the
required signs, and satisfy every exact test. Unchanged coordinates stay
fixed. The remaining inequalities hold when the choices are close enough
to the real segment, because those budgets have strict slack at \(t\).

By uniform continuity choose a finite mesh on the real segment with
successive differences less than \(\varepsilon/3\). Approximate each
interior sample within \(\varepsilon/3\) as above and use the given
rational endpoints. This gives all required successive differences.
\end{proof}

\begin{theorem}[Signed-core classification]\label{thm:spectral}
In the stated bounded class of held monomial systems, rational lawful
points \(x,y\) are finely connected if and only if \(\SCore(x)=\SCore(y)\).
\end{theorem}
\begin{proof}
For sufficiency, contract each endpoint to its core using
\Cref{lem:contraction}. Equality of signed cores puts the resulting rational
points in the same chart. Join them by \Cref{lem:rational-chart} and reverse
the chain from the second endpoint to its core.

For necessity, take the closure in \(\RR^J\) of the lawful rational
points in each chart they inhabit. There are finitely many such charts.
These closures are compact by the assumed bound and remain in the real
lawful set. If two of them meet, \Cref{lem:limit} contracts both charts
to the same signed support. By \Cref{lem:contraction}
they have the same signed core. Thus unions of chart closures for
different core values are disjoint compact sets. When there is more
than one value, their smallest pairwise distance is positive.
No rational chain of smaller mesh can cross from one such union to another,
and every lawful rational point belongs to its chart union.
This proves the necessary condition.
\end{proof}

The two uses of real analysis have separate returns. A real collision
determines a discrete support reduction. Real interpolation is then
sampled by rational points that obey the original equalities. The
reachability theorem is about finite chains at each tolerance, and in
particular does not require a continuous path taking only rational values.
